\noindent\textbf{Source:}
Folland, \emph{Real Analysis}, Second Edition.

\medskip

\noindent\textbf{Area:}
Measure Theory

\medskip

\noindent\textbf{Type:}
Foundational Problem

\medskip

\noindent\textbf{Status:}
Personal solution and structural interpretation.

\medskip

\noindent\textbf{Date:}
July 2026
\section*{Problem 01: Characterizing Sigma-Algebras by Increasing Unions}

This problem arose during my study of Measure Theory (MA222, Indian Institute of Science), following Chapter 1 of Folland's \emph{Real Analysis} (Second Edition).

Beyond presenting a complete proof, I include a structural interpretation explaining the central idea behind the argument and why this characterization appears naturally throughout modern measure theory.  

\begin{problem}

Prove that an Algebra $\mathcal{A}$ is a $\sigma$-algebra if and only if $\mathcal{A}$ is closed under increasing unions, i.e., if $\{E_i\}_{i=1}^\infty\subseteq \mathcal{A}$ and $E_1\subseteq E_2\subseteq\cdots$, then $\bigcup_{i=1}^\infty E_i\in\mathcal{A}$.  
\end{problem}
\begin{solution}
$\implies$ Assume $\mathcal{A}$ is a $\sigma$-algebra. Let $\{E_i\}_{i=1}^\infty$ be an increasing sequence i.e., $E_1\subseteq E_2\subseteq\cdots$. Since a $\sigma$-algebra is closed under arbitrary countable unions, we have:
\[
\bigcup_{i=1}^\infty E_i\in\mathcal{A}
\]
Thus, $\mathcal{A}$ is closed under increasing unions. 
Hence, the forward direction proves that if $\mathcal{A}$ is a $\sigma$-algebra, then $\mathcal{A}$ is closed under increasing unions. 

$\impliedby$
Now consider the reverse direction. Assume $\mathcal{A}$ is an algebra and closed under increasing unions. Let $\{E_n\}_{n=1}^\infty$ be an arbitrary family (not necessarily increasing). Define a new sequence by:
\[
F_n=\bigcup_{k=1}^n E_k \qquad (n\ge 1)
\]
Since $\mathcal{A}$ is an algebra, each $F_n\in\mathcal{A}$. This is trivial as an algebra is closed under finite union and $E_k\in \mathcal{A}$. Moreover $\{F_n\}$ is increasing: $F_1\subseteq F_2\cdots$. By 
\[
\bigcup_{n=1}^\infty F_n\in \mathcal{A}
\]
But $\bigcup_{n=1}^\infty F_n=\bigcup_{n=1}^\infty E_n$. Hence, $\bigcup_{n=1}^\infty E_n \in \mathcal{A}$, showing that $\mathcal{A}$ is closed under arbitrary countable unions. Now $\mathcal{A}$ already contains complements and empty set, it satisfies the axioms of $\sigma$-algebra. So, the key thing we needed to establish in the reverse direction that the given algebra is closed under arbitrary countable unions, which already proved. Now $\mathcal{A}$ satisfies all the properties of $\sigma$-algebra.\newline
Thus $\mathcal{A}$ is a $\sigma$-algebra.
\end{solution}
\begin{interpretation}
The forward direction of the above proposition is straightforward: $\mathcal{A}$ is a $\sigma$-algebra. A $\sigma$-algebra is closed under countable unions, so in particular it is closed under the special case of increasing unions. The deeper insight is required in the reverse direction. 

 Here, we start with an algebra that is closed under increasing unions. To conclude that $\mathcal{A}$ is closed under arbitrary countable unions, we transform an arbitrary family into an increasing one. The key idea is to \emph{monotonize} any given countable countable family $\{E_n\}\subset\mathcal{A}$ by considering sequence of finite unions $F_n=\bigcup_{k=1}^n E_k$. Each $F_n$ lies in $\mathcal{A}$ (since an algebra is closed under finite unions), and $F_n$ is increasing by construction. By hypothesis, $\bigcup_{n=1}^\infty F_n\in\mathcal{A}$. But $\bigcup_{n=1}^\infty F_n=\bigcup_{n=1}^\infty E_n$. This proves that $\mathcal{A}$ is closed under arbitrary countable unions.
 
 This is simple but very powerful technique---replacing an arbitrary countable unions by an increasing one, and this appears throughout measure theory, analysis. This is the essence of the monotone class theorem---an algebra that is also a monotone class (i.e., closed under increasing unions or decreasing intersections) is a $\sigma$-algebra.

 
This observation also underlines a proof that the collection of Lebesgue-measurable sets forms a $\sigma$-algebra: first by showing it is an algebra (using complements and finite unions) and then verifies closure under increasing unions via an outer measure argument. Thus the proposition is not just a technical exercise but a foundational lemma that captures the core mechanism by which $\sigma$-algebras arise in practice. 
\end{interpretation}


